\chapter{Automatic References and Bibliographies}
\label{chap:references}

\begin{learningoutcomes}
  \item Cross-reference sections, equations, tables, and figures.
  \item Build a bibliography with Bib\LaTeX{} and Biber.
  \item Keep reference data separate from prose.
\end{learningoutcomes}

\section{Label the object; refer to the label}

\begin{latexcode}{Complete example: cross-references}
\documentclass{article}
\usepackage{amsmath}
\usepackage{booktabs}
\usepackage{hyperref}
\usepackage[nameinlink,noabbrev]{cleveref}

\begin{document}
\section{Method}
\label{sec:method}

The estimate is defined in \cref{eq:estimate}.
\begin{equation}
  \hat{\theta} = \frac{1}{n}\sum_{i=1}^{n}x_i
  \label{eq:estimate}
\end{equation}

\begin{table}[htbp]
  \centering
  \caption{Simulated summary.}
  \label{tab:summary}
  \begin{tabular}{lr}
    \toprule
    Group & Mean \\
    \midrule
    A & 12.4 \\
    B & 16.7 \\
    \bottomrule
  \end{tabular}
\end{table}

\Cref{tab:summary} reports the values calculated using the
method in \cref{sec:method}.
\end{document}
\end{latexcode}

Place \cmd{label} after a section command or caption, and inside a numbered
equation. Use stable names such as \code{sec:method}, \code{eq:estimate},
\code{tab:summary}, and \code{fig:workflow}.

\section{Create a self-contained bibliography project}

The next example creates its own \file{tutorial-references.bib} so it can be typed and
compiled as one tutorial:

\begin{latexcode}{Complete example: Bib\LaTeX{} and Biber}
\begin{filecontents*}{tutorial-references.bib}
@article{measurement2026,
  author       = {Researcher, A. and Researcher, B.},
  title        = {A Reproducible Measurement Study},
  journaltitle = {Journal Name},
  year         = {2026},
  volume       = {1},
  pages        = {1--10}
}

@book{writing2025,
  author    = {Writer, C.},
  title     = {Clear Technical Writing},
  publisher = {Publisher Name},
  year      = {2025}
}
\end{filecontents*}

\documentclass{article}
\usepackage[backend=biber,style=authoryear]{biblatex}
\addbibresource{tutorial-references.bib}

\begin{document}
\textcite{measurement2026} describes the method.
Clear structure also helps readers \parencite{writing2025}.

\printbibliography
\end{document}
\end{latexcode}

Compile with:

\begin{terminalcode}{Build a bibliography project}
latexmk -pdf main.tex
\end{terminalcode}

\pkg{latexmk} runs the required \LaTeX--Biber--\LaTeX{} sequence. In a real
project, keep the records in a separate \file{references.bib} file.

\section{Useful citation commands}

\begin{latexcode}{Citation recipes}
\textcite{key}             % Author (year)
\parencite{key}            % (Author, year)
\parencite[see][p. 12]{key}
\citeauthor{key}
\citeyear{key}
\nocite{key}               % list without citing in prose
\printbibliography
\end{latexcode}

Copy bibliographic details from the source itself. Do not guess authors,
dates, titles, editions, or identifiers.

\begin{commonerror}
\textbf{Symptom:} the citation prints as a key or a question mark.

\textbf{Fix:} check that the key matches, the bibliography file is named
correctly, and Biber ran. Then compile twice more or use \pkg{latexmk}.
\end{commonerror}

\begin{codelab}
Compile the cross-reference example and move the table before the equation;
the references must remain correct. Then compile the bibliography example and
add one record of your own.
\end{codelab}

\begin{chapterchecklist}
  \item Numbers are never typed manually in cross-references.
  \item Labels are unique and use consistent prefixes.
  \item Every citation key exists in the bibliography file.
  \item Bibliographic data were checked rather than invented.
\end{chapterchecklist}
